% GENERATED FILE. Edit canonical lessons, then run npm run generate:books.
% canonical-source-hash: fnv1a64:60325e36cdb9cef7
% canonical-lessons: ZH-C50-shuijiao, ZH-C50-qichuang, ZH-C50-xiuxi, ZH-C50-shenti, ZH-C50-bing

\chapter{Sleeping and Waking}
\label{ch:zh-sleeping-and-waking}
\hlchaptermodality{50}

\begin{chapteropening}
\textbf{By the end of this chapter:} \emph{I can say when I sleep, get up and rest, and ask about health.}

Complete the last lesson of chapter 50: I can say when I sleep, get up and rest, and ask about health.
\end{chapteropening}

\section[shuìjiào]{shuìjiào — to sleep; to go to bed}

\label{lesson:ZH-C50-shuijiao}



\begin{quote}
Before the new one: say the Mandarin for afternoon, then the Mandarin for good evening.
\end{quote}

\subsection*{You'll want to know: shuìjiào}
\emph{shuìjiào} — \textquotedblleft{}to sleep; to go to bed\textquotedblright{}.

\begin{quote}
\emph{wǒ qù shuìjiào} — I'm off to bed.
\end{quote}

Its characters come later; for now it is a word for the ear.

\begin{cousinweb}
\emph{shuìjiào} is built from shuì 'to sleep' and jiào 'sleep', both inherited from Middle Chinese, \emph{dzyweH} and \emph{kaewH} (Guangyun, in Baxter's transcription); the dictionary glosses the second with this very pair.
\end{cousinweb}

\begin{grammarlens}[title={What you've built}]
One more word for saying how you are.
\end{grammarlens}

\subsection*{Your turn}
\begin{itemize}
\raggedright
  \item \emph{Say it:} \emph{shuìjiào}
  \item \emph{Say it:} \emph{shuìjiào}, once more
  \item \emph{Say it:} \emph{wǒ qù shuìjiào} — I'm off to bed.
  \item \emph{Recall:} say the Mandarin for afternoon, then the Mandarin for good evening, then say \emph{shuìjiào} again
\end{itemize}

\begin{tcolorbox}[breakable,colback=teal!4,colframe=teal!35!black,title={Before you move on}]
What does shuìjiào mean? (\textquotedblleft{}To sleep; to go to bed\textquotedblright{}.) Say it once more.
\end{tcolorbox}
\section[qǐchuáng]{qǐchuáng — to get up}

\label{lesson:ZH-C50-qichuang}



\begin{quote}
Before the new one: say the Mandarin for good evening, then the Mandarin for to sleep.
\end{quote}

\subsection*{You'll want to know: qǐchuáng}
\emph{qǐchuáng} — \textquotedblleft{}to get up\textquotedblright{}.

\begin{quote}
\emph{wǒ qī diǎn qǐchuáng} — I get up at seven.
\end{quote}

Its characters come later; for now it is a word for the ear.

\begin{cousinweb}
\emph{qǐchuáng} is built from qǐ 'to rise' and chuáng 'bed', both inherited from Middle Chinese, \emph{khiX} and \emph{dzrjang} (Guangyun, in Baxter's transcription).
\end{cousinweb}

\begin{grammarlens}[title={What you've built}]
One more word for saying how you are.
\end{grammarlens}

\subsection*{Your turn}
\begin{itemize}
\raggedright
  \item \emph{Say it:} \emph{qǐchuáng}
  \item \emph{Say it:} \emph{qǐchuáng}, once more
  \item \emph{Say it:} \emph{wǒ qī diǎn qǐchuáng} — I get up at seven.
  \item \emph{Recall:} say the Mandarin for good evening, then the Mandarin for to sleep, then say \emph{qǐchuáng} again
\end{itemize}

\begin{tcolorbox}[breakable,colback=teal!4,colframe=teal!35!black,title={Before you move on}]
What does qǐchuáng mean? (\textquotedblleft{}To get up\textquotedblright{}.) Say it once more.
\end{tcolorbox}
\section[xiūxi]{xiūxi — to rest}

\label{lesson:ZH-C50-xiuxi}



\begin{quote}
Before the new one: say the Mandarin for to sleep, then the Mandarin for to get up.
\end{quote}

\subsection*{You'll want to know: xiūxi}
\emph{xiūxi} — \textquotedblleft{}to rest\textquotedblright{}.

\begin{quote}
\emph{xīngqītiān wǒ xiūxi} — On Sunday I take it easy.
\end{quote}

Its characters come later; for now it is a word for the ear.

\begin{cousinweb}
\emph{xiūxi} is built from xiū 'to rest' and xi 'to stop, to breathe', both inherited from Middle Chinese, \emph{xjuw} and \emph{sik} (Guangyun, in Baxter's transcription).
\end{cousinweb}

\begin{grammarlens}[title={What you've built}]
One more word for saying how you are.
\end{grammarlens}

\subsection*{Your turn}
\begin{itemize}
\raggedright
  \item \emph{Say it:} \emph{xiūxi}
  \item \emph{Say it:} \emph{xiūxi}, once more
  \item \emph{Say it:} \emph{xīngqītiān wǒ xiūxi} — On Sunday I take it easy.
  \item \emph{Recall:} say the Mandarin for to sleep, then the Mandarin for to get up, then say \emph{xiūxi} again
\end{itemize}

\begin{tcolorbox}[breakable,colback=teal!4,colframe=teal!35!black,title={Before you move on}]
What does xiūxi mean? (\textquotedblleft{}To rest\textquotedblright{}.) Say it once more.
\end{tcolorbox}
\section[shēntǐ]{shēntǐ — body; health}

\label{lesson:ZH-C50-shenti}



\begin{quote}
Before the new one: say the Mandarin for to get up, then the Mandarin for to rest.
\end{quote}

\subsection*{You'll want to know: shēntǐ}
\emph{shēntǐ} — \textquotedblleft{}body; health\textquotedblright{}.

\begin{quote}
\emph{nǐ shēntǐ hǎo ma} — How's your health?
\end{quote}

Its characters come later; for now it is a word for the ear.

\begin{cousinweb}
\emph{shēntǐ} is built from shēn 'body, self' and tǐ 'body', both inherited from Middle Chinese, \emph{syin} and \emph{thejX} (Guangyun, in Baxter's transcription).
\end{cousinweb}

\begin{grammarlens}[title={What you've built}]
One more word for saying how you are.
\end{grammarlens}

\subsection*{Your turn}
\begin{itemize}
\raggedright
  \item \emph{Say it:} \emph{shēntǐ}
  \item \emph{Say it:} \emph{shēntǐ}, once more
  \item \emph{Say it:} \emph{nǐ shēntǐ hǎo ma} — How's your health?
  \item \emph{Recall:} say the Mandarin for to get up, then the Mandarin for to rest, then say \emph{shēntǐ} again
\end{itemize}

\begin{tcolorbox}[breakable,colback=teal!4,colframe=teal!35!black,title={Before you move on}]
What does shēntǐ mean? (\textquotedblleft{}Body; health\textquotedblright{}.) Say it once more.
\end{tcolorbox}
\section[bìng]{bìng — to fall ill; illness}

\label{lesson:ZH-C50-bing}



\begin{quote}
Before the new one: say the Mandarin for to rest, then the Mandarin for health.
\end{quote}

\subsection*{You'll want to know: bìng}
\emph{bìng} — \textquotedblleft{}to fall ill; illness\textquotedblright{}.

\begin{quote}
\emph{wǒ māma bìng le} — My mum is ill.
\end{quote}

Its characters come later; for now it is a word for the ear.

\begin{cousinweb}
\emph{bìng} is inherited from Middle Chinese \emph{bjaengH} (Guangyun, in Baxter's transcription), glossed 'worry, pain' and 'an illness grown worse'.
\end{cousinweb}

\begin{grammarlens}[title={What you've built}]
One more word for saying how you are.
\end{grammarlens}

\subsection*{Your turn}
\begin{itemize}
\raggedright
  \item \emph{Say it:} \emph{bìng}
  \item \emph{Say it:} \emph{bìng}, once more
  \item \emph{Say it:} \emph{wǒ māma bìng le} — My mum is ill.
  \item \emph{Recall:} say the Mandarin for to rest, then the Mandarin for health, then say \emph{bìng} again
\end{itemize}

\begin{tcolorbox}[breakable,colback=teal!4,colframe=teal!35!black,title={Before you move on}]
What does bìng mean? (\textquotedblleft{}To fall ill; illness\textquotedblright{}.) Say it once more.
\end{tcolorbox}
